\documentclass[10pt]{article}
\usepackage[margin=0.78in]{geometry}
\usepackage[T1]{fontenc}
\usepackage{lmodern}
\usepackage{booktabs}
\usepackage{longtable}
\usepackage{graphicx}
\usepackage{xcolor}
\usepackage{hyperref}
\usepackage{fancyhdr}
\usepackage{microtype}
\hypersetup{colorlinks=true,linkcolor=blue!55!black,urlcolor=blue!55!black}
\definecolor{stopred}{HTML}{9A2D2D}
\definecolor{auditblue}{HTML}{2C6E9E}
\pagestyle{fancy}
\fancyhf{}
\lhead{PD39 physical blocker validation}
\rhead{CASE C - STOP}
\cfoot{\thepage}
\setlength{\parindent}{0pt}
\setlength{\parskip}{5pt}
\newcommand{\code}[1]{\texttt{#1}}
\newcommand{\fail}{\textcolor{stopred}{FAIL}}
\title{PD39 Physical Blocker Validation\\\large First Compatibility Atlas - STOP Report}
\author{PD39 research branch: \code{research/pd39-physical-blockers-atlas-v1}}
\date{14 September 2026}

\begin{document}
\maketitle

\begin{center}
{\Large \textcolor{stopred}{FINAL CASE: C - STOP}}\\[4pt]
No physically confirmed blocker; no compatibility atlas was run.
\end{center}

\section*{Executive verdict}

The prior fresh-holdout census supplied three low-order C12 portfolios that
looked like H0 candidates on the default PowerDynamics initialization path:
\code{35;36}, \code{37;38}, and \code{30;32;33}. This campaign tested those
cases, their immediate predecessors, cheap proper subsets, and aggregate-matched
controls. The default-path positive alphas were not reproducible under the
frozen Gate A tolerance/path audit. All three selected C12 blocker TDS runs
returned \code{MaxIters}; predecessors and controls completed and decayed.

Thus the candidates are not physically validated dynamic instabilities. The
correct result is a model/equilibrium-path warning. The conditional atlas,
composition claim, contextual-return validation, remediation, and excluded
campaigns were not run.

\section{Scope, provenance, and frozen decision}

The final branch is \code{research/pd39-physical-blockers-atlas-v1}, created
from parent commit \code{e81d18ab945b}. Preregistration was committed before
new numerical results at \code{1b9ad5fe}. Gate A, default-path audit, and Gate
B commits were \code{02381549}, \code{d24f0ad3}, and \code{c95bab91}.
PowerDynamics was version 5.0.0 under Julia 1.11.9. No push was performed.

Only the following gates were allowed: numerical/DAE physicality, common
nonlinear TDS, conditional aggregate composition, and exact network-closure
feasibility. Weak-node/link ranking, structured-radius search, repair
optimization, planners, co-design, IEEE-68, EMT, and new controller search
were explicitly out of scope.

The frozen minimality separation was
\[
 \delta_H = \min\left(\alpha(H), -\max_{G\subset H}\alpha(G)\right),
\]
with primary threshold \code{0.02 s-1}. A positive alpha alone was not
accepted as physicality evidence.

\section{Discovery audit and frozen selection}

The authoritative prior file contained 256 portfolios x 24 holdout conditions,
or 6144 portfolio-condition evaluations. Its frozen summary reported 42
minimal H0 cases and 47 minimal H0.05 cases. This narrow campaign did not
rerun the census.

\begin{table}[h]
\centering
\caption{Frozen candidate selection.}
\begin{tabular}{llllr}
\toprule
ID & Candidate & Condition & Matched control & Distance\\
\midrule
H1 & 35;36 & C12 & 32;36 & 0.0013860\\
H2 & 37;38 & C12 & 36;38 & 0.0148840\\
H3 & 30;32;33 & C12 & 30;33;35 & 0.0014786\\
\bottomrule
\end{tabular}
\end{table}

The controls were frozen by same cardinality, then MW, MVA, remaining
inertia, and lexicographic order. No selection was changed after observing
results.

\section{Gate A - numerical and DAE physicality}

The audit covered 17 unique portfolios at C12 and nominal with tolerances
\code{1e-8}, \code{1e-10}, and \code{1e-12}. It used the native reduction,
an independent dense spectrum, a descriptor/generalized check where exposed,
central finite-difference diagnostics, conditioning, singular values, and
participation. Within each fixed solved path, dense and descriptor spectra
agreed with the native result. The equilibrium path did not.

\begin{table}[h]
\centering
\caption{Default census initialization path at C12.}
\begin{tabular}{lrrll}
\toprule
ID & alpha & frequency (Hz) & mode label & residual\\
\midrule
H1 & +10.1875267 & 0.000000 & electromechanical/control & 3.54e-12\\
H2 & +22.5280341 & 0.549890 & converter-control oscillatory & 7.04e-13\\
H3 & +21.1535104 & 0.000000 & converter-control oscillatory & 7.04e-13\\
\bottomrule
\end{tabular}
\end{table}

\begin{table}[h]
\centering
\caption{Strict Gate A path at tolerance 1e-10.}
\begin{tabular}{lrrl}
\toprule
ID & alpha & frequency (Hz) & interpretation\\
\midrule
H1 & -0.0979116 & 0.067999 & stable on this path\\
H2 & +11.8934740 & 0.000000 & positive but path-sensitive\\
H3 & -0.1088351 & 0.074147 & stable on this path\\
\bottomrule
\end{tabular}
\end{table}

At the default path, the immediate-predecessor separations for H1, H2, and H3
were 0.0971024, 0.0975787, and 0.0978716 s$^{-1}$. Those values are not
physical minimality evidence once the candidate branch is shown to be
path-sensitive. Large reduced-Jacobian conditioning, near-zero singular
values, and PowerDynamics vertex-batch warnings are consistent with a
numerical/DAE branch concern, but do not select a physical mechanism.

\textbf{Gate A: \fail\_UNRESOLVED.}

\begin{figure}[h]
\centering
\includegraphics[width=0.98\linewidth]{../figures/PD39_blocker_atlas/F1_path_dependence_selected_cases.pdf}
\caption{Default versus strict audit path for candidates and controls.}
\end{figure}

\section{Gate B - common nonlinear TDS}

The preregistered disturbance was bus 39 with a +1\% P/Q constant-power-factor
pulse from 1.0 to 1.1 s, restoration, and integration to 20 s with common
Rodas5P settings. The run contained 26 rows. A \code{maxiters=100000} guard
was added only to prevent unbounded runtime; it did not change the disturbance,
tolerances, observables, or acceptance rule. A max-iteration return is a TDS
failure.

All three selected C12 candidates returned \code{MaxIters}. Immediate
predecessors and controls completed with negative estimated rates, roughly
-0.207 to -0.258 s$^{-1}$. The failed candidate integrations showed large
state/Jacobian excursions before the solver stopped. This pattern is not
nonlinear physical confirmation.

\textbf{Gate B: \fail.}

\begin{figure}[h]
\centering
\includegraphics[width=0.98\linewidth]{../figures/PD39_blocker_atlas/F2_tds_stop_summary.pdf}
\caption{C12 TDS completion pattern. Red candidate bars are absent because
the solver returned MaxIters; blue neighbors completed and decayed.}
\end{figure}

\section{Gate C - composition}

The composition gate was conditional on physical blocker confirmation. Because
Gates A and B did not confirm a blocker, the gate was not entered. Raw
default-path candidate/control comparisons are retained in
\code{results/PD39\_BLOCKER\_COMPOSITION\_TEST.csv}; they do not constitute a
composition claim. In particular, a large H2 default-path contrast cannot be
promoted while the equilibrium branch is unresolved.

\textbf{Gate C: NOT APPLICABLE; composition strong = false.}

\section{Penetration, composition, and mode mechanism}

The prior holdout results show H0/H0.05 cases across multiple conditions rather
than a universal V8 identity. In the selected C12 cases, the default labels
are heterogeneous: H1 is electromechanical/control and H2/H3 are labeled
converter-control by the default path. The tolerance/path changes and failed
TDS integrations prevent a physical 7/8 -> 8/8 mechanism claim. No physical
SG-to-GFL homotopy was run.

\begin{figure}[h]
\centering
\includegraphics[width=0.98\linewidth]{../figures/PD39_blocker_atlas/F3_gate_a_tolerance_sweep.pdf}
\caption{Frozen tolerance sweep for selected C12 candidates.}
\end{figure}

\section{Exact network closure}

The installed model/API does not expose the exact device-port objects, local
admittance partition, and dimensioned K, D, and Q objects required for the
requested contextual-return identity. PD contextual return is therefore
\code{BLOCKED}. TX4 contextual-return validation was not run after the
physicality STOP. No proxy, determinant identity, or Q -> -1 boundary was
constructed.

\textbf{Closure: BLOCKED.}

\section{Compatibility atlas and conditional theory}

The PLL/filter two-coordinate atlas, 11 x 11 grid, adaptive boundaries,
\code{kappa\_0}, \code{kappa\_0.05}, mode-MAC continuation, contextual-return
theorem validation, and local remediation were not run. A compatibility atlas
on path-sensitive equilibria and failed TDS cases would overstate the evidence.

\section{Negative results and forbidden claims}

This branch does not support a physical minimal H0/H0.05 blocker, a physical
7/8 -> 8/8 mode transition, an aggregate composition mechanism, compatibility
boundaries, contextual-return closure, or any controller/line/topology repair.
It also does not revive hidden-radius, generic weak-node/link, graph
Laplacian/Fiedler, cycle/holonomy, Shapley, or broad co-design novelty.

The strongest defensible positive statement is narrower: the stock PD39
workflow produces low-order C12 instability candidates whose default-path
signs are not reproducible under the frozen audit and whose common-disturbance
trajectories can terminate in solver iteration failure while neighbors decay.
This is a validation warning requiring model/equilibrium-path repair.

\section{Practical interpretation and final decision}

For SG-to-IBR planning, the immediate lesson is a quality-control rule: do not
promote a portfolio to a physical blocker from one initialization and one
rightmost eigenvalue. Require branch reproducibility, independent spectrum
agreement, conditioning diagnostics, and a common nonlinear disturbance that
completes.

The campaign does not justify redesigning the IAS poster around a physical
blocker or atlas. It is not Transactions-ready as a new mechanism result.

\begin{table}[h]
\centering
\caption{Final status matrix.}
\begin{tabular}{ll}
\toprule
Item & Status\\
\midrule
Numerical audit & FAIL\_UNRESOLVED\\
Selected-blocker TDS & FAIL: 3 C12 MaxIters\\
Prior H0 count over 24 holdouts & 42\\
Prior H0.05 count over 24 holdouts & 47\\
Penetration/composition/mixed verdict & no physical blocker\\
Mode mechanism & unresolved, path-dependent labels\\
Closure & BLOCKED\\
Atlas & NOT RUN\_STOP\\
IAS poster redesign & NO\\
Transactions-ready & NO\\
\bottomrule
\end{tabular}
\end{table}

\textbf{FINAL CASE: C.}

\section*{Reproducibility files}

The final repository retains the preregistration, deviations, selection table,
Gate A tables, default-path audit, Gate B raw traces, STOP figures, final
machine-readable claims, and two ZIP bundles. The compact bundle is
\code{deliverables/PD39\_BLOCKER\_ATLAS\_CHATGPT\_UPLOAD.zip}.

\end{document}
